% =====================================================================
% School Wits — QUESTION PAPER (QP) template
%
% This is the canonical preamble shared by every QP file in
% "Updated Latex FIles For Web/". Copy this file, keep the preamble
% BYTE-IDENTICAL, and write only inside \begin{document}...\end{document}.
%
% The preamble is not decoration — the parser keys off these exact macro
% names to turn the paper into JSON. Renaming \examq or \qfig, or
% hand-rolling a question heading with \section instead, means that
% question will not be extracted.
%
% See README.md in this folder for the macro reference and the JSON
% shape each macro maps to.
% =====================================================================

\documentclass[11pt,a4paper]{article}

\usepackage[T1]{fontenc}
\usepackage[utf8]{inputenc}
\usepackage{lmodern}
\usepackage[a4paper,top=2.4cm,bottom=2.2cm,left=2.2cm,right=2.2cm]{geometry}
\usepackage{amsmath}
\usepackage{amssymb}
\usepackage{graphicx}
\graphicspath{{./}}
\usepackage{array}
\usepackage{makecell}
\renewcommand{\arraystretch}{1.25}
\usepackage{enumitem}
\usepackage{fancyhdr}
\usepackage{xcolor}

% ---- Subject-specific shorthands -------------------------------------
% These four are all any QUESTION paper uses — the Add Maths shorthands
% (\dd, \dydx, \ncr, \npr, \cosec) appear only in mark schemes, so they
% live in answer-paper.template.tex instead. Anything you add here must
% also be taught to KaTeX in js/app.js's KATEX_MACROS, or it will render
% as raw source in the browser.
\newcommand{\ohms}{\ensuremath{\Omega}}
\newcommand{\degC}{\ensuremath{^{\circ}}C}
\newcommand{\dg}{\ensuremath{^{\circ}}}
\newcommand{\pow}[1]{\ensuremath{^{#1}}}

\newcommand{\topiclabel}[1]{{\hfill\normalfont\small\itshape\textcolor{black!55}{#1}}}

% ---- Answer-space macros (layout only, no content) -------------------
% These exist so the printed PDF has writing space. The parser treats
% them as "blank answer space" and uses only the mark value in the last
% argument, e.g. \Alines{3}{[1]} contributes marks "1".
\newlength{\anslineskip}
\setlength{\anslineskip}{8mm}

\newcounter{ansln}
\newcommand{\Alines}[2]{%
  \par\vspace{1mm}%
  \setcounter{ansln}{1}%
  \loop\ifnum\value{ansln}<#1\relax
    \noindent\mbox{}\dotfill\par\vspace{\anslineskip}%
    \stepcounter{ansln}%
  \repeat
  \noindent\mbox{}\dotfill\hspace{0.6em}\mbox{#2}\par
  \vspace{1mm}%
}

\newcommand{\plainline}{\par\noindent\mbox{}\dotfill\par\vspace{\anslineskip}}

\newcommand{\labelline}[1]{\par\vspace{1mm}\noindent #1\hspace{0.5em}\dotfill\par\vspace{\anslineskip}}

% \ansval{label}{unit}{[marks]} — a right-aligned "x = ....... [2]" line.
\newcommand{\ansval}[3]{%
  \par\vspace{5mm}%
  \noindent\hspace*{3.5cm}#1\hspace{0.5em}\dotfill\hspace{0.5em}#2\hspace{0.5em}\makebox[1.1cm][r]{#3}\par
  \vspace{3mm}%
}

% \markright{[2]} — a right-aligned mark tag for a part with no ruled
% answer space (e.g. one answered by drawing on a figure).
\renewcommand{\markright}[1]{\par\nopagebreak\noindent\hspace*{\fill}\mbox{#1}\par\vspace{1mm}}

% ---- THE ONE MACRO THAT MATTERS --------------------------------------
% \examq{paper-id}{number}{topics}{marks}
%   paper-id : syllabus code + variant + session + year, e.g. 5054/21/M/J/25
%   number   : question number, bare integer, no parentheses
%   topics   : one or more topic names joined by \textperiodcentered\
%   marks    : total marks for the question, bare integer
% This is the ONLY question delimiter. Everything between one \examq and
% the next belongs to that question.
\newcommand{\swmarks}{}
\newcommand{\examq}[4]{%
  \gdef\swmarks{#4}%
  \par\addvspace{16pt}\noindent
  {\large\bfseries #1 - Q#2 - #3 - #4}\par\vspace{5pt}%
}

% \total takes NO argument: it prints the marks stored by \examq.
% Put it as the last line of every structured question.
\newcommand{\total}{\par\vspace{2mm}\noindent\hspace*{\fill}[Total:\ \swmarks]\par\vspace{2mm}}

% \qfig[width]{filename}{caption} — caption may be empty: {}
% The filename must exist next to the .tex file.
\newcommand{\qfig}[3][0.5\linewidth]{%
  \par\vspace{2mm}
  \begin{center}%
    \IfFileExists{#2}%
      {\includegraphics[width=#1]{#2}}%
      {\fbox{\parbox[c][26mm][c]{#1}{\centering\ttfamily\footnotesize #2}}}\\[4pt]
    \textbf{#3}%
  \end{center}%
  \vspace{2mm}\par
}

% ---- Part numbering --------------------------------------------------
% parts    -> (a) (b) (c)     subparts -> (i) (ii) (iii)
% subparts may only be nested INSIDE a parts \item.
\newlist{parts}{enumerate}{1}
\setlist[parts]{label=\textbf{(\alph*)},leftmargin=2.1em,labelsep=0.6em,itemsep=10pt,topsep=6pt,parsep=4pt}

\newlist{subparts}{enumerate}{1}
\setlist[subparts]{label=\textbf{(\roman*)},leftmargin=2.3em,labelsep=0.6em,itemsep=10pt,topsep=6pt,parsep=4pt}

% ---- MCQ papers only -------------------------------------------------
% Multiple-choice option list -> A B C D.
% Use this for EVERY multiple-choice question. Do not hand-roll options
% with \textbf{A}\hspace{...} or a bare tabular — see README.md, issue #1.
\newlist{choices}{enumerate}{1}
\setlist[choices]{label=\textbf{\Alph*},leftmargin=2.1em,labelsep=0.9em,itemsep=2pt,topsep=5pt,parsep=0pt}

% ---- Running head: update these three lines per paper ----------------
\pagestyle{fancy}
\fancyhf{}
\fancyhead[L]{\footnotesize 5054/21/M/J/25}
\fancyhead[C]{\footnotesize Cambridge O Level Physics -- Paper 2}
\fancyhead[R]{\footnotesize May/June 2025}
\fancyfoot[C]{\footnotesize\thepage}
\renewcommand{\headrulewidth}{0.4pt}
\renewcommand{\footrulewidth}{0pt}

\setlength{\parindent}{0pt}

\begin{document}

%============================== QUESTION 1 ==============================
% A STRUCTURED question: stem -> figure -> parts -> subparts -> \total
\examq{5054/21/M/J/25}{1}{Motion or Kinematics \textperiodcentered\ Forces or Dynamics}{9}
Fig. 1.1 shows a skydiver falling vertically through the air.

\qfig[0.42\linewidth]{5054_21_M_J_25_fig1.png}{Fig. 1.1}

In the first part of the fall, her speed increases and her acceleration decreases.

\begin{parts}
  % (a) — answered by drawing on a figure, so \markright supplies the mark
  \item On Fig. 1.2 sketch the speed--time graph for the skydiver.

  \qfig[0.55\linewidth]{5054_21_M_J_25_fig2.png}{Fig. 1.2}
  \markright{[2]}

  % (b) — ruled answer space: 3 lines, worth 1 mark
  \item Explain how the graph shows that the acceleration decreases as the speed increases.
  \Alines{3}{[1]}

  % (c) — a part that contains subparts (i), (ii)
  \item During the first part of the fall, there is a resultant vertical force
  acting downwards on the skydiver.
  \begin{subparts}
    \item State the name of the other vertical force that acts on the skydiver.
    \Alines{1}{[1]}

    \item Explain why the resultant vertical force eventually becomes zero.
    \Alines{4}{[2]}
  \end{subparts}

  % (d) — a numeric answer with a labelled value line
  \item At one instant, the vertical force on the skydiver is 400\,N downwards.

  Determine the resultant force.
  \ansval{resultant $=$}{N}{[3]}
\end{parts}
\total

%============================== QUESTION 2 ==============================
% An MCQ question. Note: NO \begin{parts}, NO \total — marks are always 1.
\examq{5054/11/M/J/25}{2}{Momentum (Dynamics)}{1}
What is a unit for momentum?

\begin{choices}
  \item kg\,m\,/\,s
  \item kg\,m\,/\,s\pow{2}
  \item N\,m
  \item N\,/\,m\pow{2}
\end{choices}

\end{document}
